
\section{Data Collection}
\label{sec:dataset}



\begin{figure}[t]
  \centering
  \includegraphics[width=\linewidth]{figures/fig-methodology-3.pdf}
  \caption{Methodology-figure}
  \label{fig:figure2}
\end{figure}



\begin{table}[h]
\centering
\caption{Number of projects retained at each project-selection step. PS7 and PS8 were collected only for JavaScript projects.}
\label{tab:project-selection}
\resizebox{\columnwidth}{!}{%
\begin{tabular}{lrrrrrrrr}
\hline
 % &  & & & & & RQ2 \& RQ3.2 & & RQ1 \& RQ3.1  \\
Language & Collected & PS1 & PS2 & PS3  & PS4  & \textbf{PS5}  & PS6  & \textbf{PS7}  \\
\hline
JavaScript & 294,477 & 46,249 & 22,716 & 18,286 & 14,705 & 9,330    & 6,322 & 260 \\
Python     & 122,399 & 24,940 & 9,018  & 6,664  & 4,327  & 3,184    & --  & --  \\
Rust       & 15,300  & 5,394  & 3,702  & 2,724  & 2,539  & 2,163    & --   & --  \\
C\#        & 26,800  & 4,263  & 1,861  & 1,316  & 1,258  & 934    & --   & --  \\
\\
Total & 458, 976 & & & & &  15, 611 & & \\
\hline
\end{tabular}%
}
\end{table}

In order to answer the research questions, we collect a dataset of 458,976 publicly available GitHub repositories across four major languages (JavaScript, Python, Rust, and C\#). Using the search endpoints of the GitHub REST API\footnote{\url{https://docs.github.com/en/rest/search/search?apiVersion=2026-03-10}}, we extracted all repositories with at least 10 stars in each of the languages. This process resulted in 294,477 JavaScript repositories, 122,399 Python repositories, 15,300 Rust repositories, and 26,800 C\# repositories. Following this, we then filtered the initial dataset for projects that had dependencies, active, with recent CI use and were of a reasonable size, to ensure projects were representative of real projects in use. \autoref{tab:project-selection} presents the results of project selection. The details of the filtering are described below.

\begin{enumerate}[label={\textbf{(PS\arabic*).}}, leftmargin=50pt]
    \item \textbf{Recent activity.}
    We retained only repositories with at least 10 commits during the study period from March 1, 2024, to March 1, 2026, thereby restricting the dataset to projects that showed continuous development activity. This process resulted in 46,249 JavaScript repositories, 24,940 Python repositories, 5,394 Rust repositories, and 4,263 C\# repositories.
    
    \item \textbf{Using GitHub Actions.}
    We selected repositories that contained workflow files with either the \texttt{.yml} or \texttt{.yaml} extension under the \texttt{.github/workflows/} directory, that is, projects adopting GitHub Actions. This process resulted in 22,716 JavaScript repositories, 9,018 Python repositories, 3,702 Rust repositories, and 1,861 C\# repositories.
    
    \item \textbf{Active CI Records.}
    By examining the execution history of GitHub Actions through the GitHub API, we retained only repositories with at least 10 recorded workflow runs. This step ensured that CI was actively used in practice. This process resulted in 18,286 JavaScript repositories, 6,664 Python repositories, 2,724 Rust repositories, and 1,316 C\# repositories.
    
    \item \textbf{Has dependencies.}
    Across languages, projects use ecosystem specific configuration files to record their dependencies. Using the Contents endpoint of the GitHub REST API\footnote{\url{https://docs.github.com/en/rest/repos/contents?apiVersion=2026-03-10}}, we checked for the presence of ecosystem-specific configuration files in each repository and retained only projects that adopted the standard package management or build system for their primary language. Specifically, we selected JavaScript projects containing \texttt{package.json}, Python projects containing \texttt{pyproject.toml}, Rust projects containing \texttt{Cargo.toml}, and C\# projects containing \texttt{.csproj}. This step ensured that the selected repositories were actual software projects developed within the standard ecosystem of each language, and that they made use of dependencies. This process resulted in 14,705 JavaScript repositories, 4,327 Python repositories, 2,539 Rust repositories, and 1,258 C\# repositories.
    \item \textbf{Repository size.}
    Due to the limitations of our approach, in order to ensure that the source files of each repository can fit within the context window of the LLMs used in our analysis, we retained only repositories less than 10 MB (as reported by the GitHub API size field\footnote{\url{https://docs.github.com/en/rest/repos/repos?apiVersion=2026-03-10}}). This process resulted in 9,330 JavaScript repositories, 3,184 Python repositories, 2,163 Rust repositories, and 934 C\# repositories.\\
\end{enumerate}

The result of our filtering is a dataset of 15,611 repositories across four languages that represent active projects within their respective ecosystems.

\subsection{Testable Projects (RQ1)}
In order to answer evaluate the accuracy of the unused dependency detection approaches in RQ1, we measure weather a removing a dependency predicted as unused preserved successful test execution. However, in order to do this, we require a subset of projects with executable tests. For this subset, we apply the additional filters PS6 and PS7, resulting in a dataset of 260 Javascript repositories.


\begin{enumerate}[label={\textbf{(PS\arabic*).}}, leftmargin=50pt]
\setcounter{enumi}{5}
    \item \textbf{Verification of Test Script Availability.}
    For each language, we retrieved the relevant configuration file via the GitHub REST API and retained only repositories in which a test framework or test script was explicitly configured.
    For JavaScript, we retrieved the \texttt{package.json} file via the GitHub Contents API and retained only repositories that defined a \texttt{test} entry in the \texttt{scripts} field, resulting in 6,322 repositories.
    For Python, we retrieved the \texttt{pyproject.toml} file via the GitHub Contents API and retained only repositories in which \texttt{pytest} was declared as a dependency (e.g., within \texttt{[dependency-groups]} or \texttt{[project.optional-dependencies]}), resulting in 1,963 repositories.
    For Rust, we retrieved the complete file tree of each repository via the GitHub Tree API and retained only repositories containing a \texttt{tests/} or \texttt{test/} directory at any level of the tree hierarchy, indicating the presence of integration tests.
    For C\#, we retrieved all \texttt{.csproj} files via the GitHub Tree API and Contents API and retained only repositories in which at least one project file referenced the xUnit testing framework.
    \item \textbf{Verification of Test Executability and Coverage.}
    We performed a shallow clone (\texttt{git clone --depth=1}) of each repository, ran \texttt{npm install}, and measured test coverage using \texttt{nyc}. Only projects whose tests passed and whose line coverage was at least 70\% were retained. This process resulted in 260 JavaScript repositories.
\end{enumerate}

The dataset comprises \textbf{260} projects and includes \textbf{1,190} examined runtime dependencies and \textbf{3,234} examined development dependencies.
Commit counts range from 11 to 6,131 (median: 286.5), and GitHub star counts
range from 10 to 68,804 (median: 66), reflecting a mix of small utility
libraries and widely adopted packages.

Regarding test coverage, all projects satisfy the PS8 minimum of 70\%
nyc line coverage (median: 97.4\%).
The coverage distribution was as follows: 21 projects fell in the range $[70, 80)$\%, 35 in $[80, 90)$\%, 116 in $[90, 100)$\%, and 88 achieved perfect 100\% coverage.

To the best of our knowledge, the dataset produced by PS7 is the first publicly available benchmark for evaluating unused dependency detection in npm projects under the constraint of executable, high-coverage test suites. We make this dataset publicly available on Zenodo: XXXXXX.